\section{*Squeezed states of light (11/4)}

\referto{\scully~2.5-2.7}

\subsection{Squeezed states}

\subsubsection{Field quadratures}
For a single mode, define the field quadratures:
\begin{equation}
\hat X=\frac{\hat a+\hat a^\dagger}{\sqrt2},
\qquad
\hat Y=\frac{\imth(\hat a^\dagger-\hat a)}{\sqrt2}.
\end{equation}
Then
\begin{equation}
\comm{\hat X}{\hat Y}=\imth,
\qquad
\Delta X\,\Delta Y\geq\frac12.
\end{equation}
For coherent states,
\begin{equation}
\Delta X=\Delta Y=\frac{1}{\sqrt2}.
\end{equation}

\subsubsection{Squeezing}
A minimum-uncertainty squeezed state satisfies
\begin{equation}
\left(\Delta\hat X+\imth R^2\Delta\hat Y\right)\dket{\psi}=0.
\end{equation}
Equivalently,
\begin{equation}
\left[(1-R^2)\hat a^\dagger+(1+R^2)\hat a
-(X_0+\imth R^2Y_0)\right]\dket{\psi}=0.
\end{equation}
Define the squeeze parameter by
\begin{equation}
\frac{1-R^2}{1+R^2}=\tanh r,
\qquad
R=e^{-r}.
\end{equation}
Then
\begin{equation}
\boxed{\Delta X=\frac{e^{-r}}{\sqrt2},
\qquad
\Delta Y=\frac{e^r}{\sqrt2},
\qquad
\Delta X\Delta Y=\frac12.}
\end{equation}

\subsubsection{Electric-field noise}
For phase
\begin{equation}
\Phi=\omega t-\bm k\cdot\bm r,
\end{equation}
the electric field (\ref{electric field}) is proportional to
\begin{equation}
\hat E(\bm r,t)=\sqrt{\frac{2\hbar\omega}{\epsilon V}}
\left(\hat X\cos\Phi+\hat Y\sin\Phi\right).
\end{equation}
For a squeezed state with vanishing quadrature covariance,
\begin{equation}
\boxed{
(\Delta E)^2=\frac{\hbar\omega}{\epsilon V}
\left(e^{-2r}\cos^2\Phi+e^{2r}\sin^2\Phi\right).}
\end{equation}
For $r>0$,
\begin{align}
\Delta E_{\min}&=\Delta X=\frac{e^{-r}}{\sqrt2},
&&\Phi=0\pmod\pi,\\
\Delta E_{\max}&=\Delta Y=\frac{e^r}{\sqrt2},
&&\Phi=\frac\pi2\pmod\pi,\\
\Delta E_{\min}\Delta E_{\max}&=\Delta X\Delta Y=\frac12.
\end{align}

\subsection{Generation of coherent and squeezed states in quantum optics}

\subsubsection{Generating coherent states}
With a classical current,
\begin{equation}
\hat H=\hat H_{\mathrm{EM}}
-\int\bm j(\bm r,t)\cdot\hat{\bm A}(\bm r)\text{d}^3{\bm r}.
\end{equation}
The evolved vacuum is a coherent state:
\begin{equation}
e^{-\imth\hat Ht/\hbar}\dket{0}
=\exp\left[\sum_k(\alpha_k\hat a_k^\dagger
-\alpha_k^*\hat a_k)\right]\dket{0},
\end{equation}
where
\begin{equation}
\alpha_k=-\int\text{d}^3{\bm r}\int_0^t\dif t'
\,e^{\imth(\omega t'-\bm k\cdot\bm r)}
\bm j(\bm r,t')\cdot\bm e_k
\sqrt{\frac{\hbar}{\omega\epsilon V}}.
\end{equation}
The displacement operator is
\begin{equation}
\boxed{\hat D(\alpha)=e^{\alpha\hat a^\dagger-\alpha^*\hat a}}.
\end{equation}

\subsubsection{Generating squeezed states}
Degenerate parametric amplification, achieved by placing a nonlinear crystal nonlinear inside an optical cavity, produces the effective Hamiltonian
\begin{equation}
\hat H_{\mathrm{dpa}}
=\hbar\chi\beta_0
\left[(\hat a^\dagger)^2e^{-\imth\phi}
+\hat a^2e^{\imth\phi}\right].
\end{equation}
Define the squeeze operator:
\begin{equation}
\boxed{
\hat S(\xi)=
\exp\left[\frac12\left(\xi^*\hat a^2
-\xi(\hat a^\dagger)^2\right)\right],
\qquad \xi=re^{\imth\theta}.}
\end{equation}
It satisfies
\begin{equation}
\hat S^\dagger(\xi)(\hat X+\imth\hat Y)\hat S(\xi)
=e^{-r}\hat X+\imth e^r\hat Y
\end{equation}
for a suitable squeezing phase. A squeezed coherent state is
\begin{equation}
\boxed{
\dket{\xi,z}=\hat S(\xi)\hat D(z)\dket{0}
\sim\hat S(\xi)\dket{z}.}
\end{equation}
Its quadrature variances are
\begin{equation}
\dbra{\xi,z}(\Delta\hat X)^2\dket{\xi,z}
=\frac{e^{-2r}}{2},
\qquad
\dbra{\xi,z}(\Delta\hat Y)^2\dket{\xi,z}
=\frac{e^{2r}}{2}.
\end{equation}


\subsection{Application: squeezed light in LIGO}

\subsubsection{Gravitational-wave measurement}

LIGO is a Michelson interferometer with long optical cavities in its two
perpendicular arms. A passing gravitational wave changes the arm lengths in
opposite directions. If the unperturbed arm length is \(L\), define the
gravitational-wave strain by
\begin{equation}
    h(t)=\frac{\Delta L_x(t)-\Delta L_y(t)}{L}.
\end{equation}
The differential change in arm length produces a phase shift in the light
returning from the two arms. The output photodetector measures this phase
shift.

Ideally, the interferometer is operated near a dark fringe, so little
classical carrier light exits through the antisymmetric port. Quantum vacuum
fluctuations can nevertheless enter through this port and interfere with the
carrier field. These fluctuations contribute quantum noise to the measured
signal.

\subsubsection{Two sources of quantum noise}

Introduce the amplitude and phase quadratures
\begin{equation}
    \hat X=\frac{\hat a+\hat a^\dagger}{\sqrt{2}},
    \qquad
    \hat Y=\frac{\imth(\hat a^\dagger-\hat a)}{\sqrt{2}},
\end{equation}
with
\begin{equation}
    \comm{\hat X}{\hat Y}=\imth,
    \qquad
    \Delta X\,\Delta Y\geq\frac12.
\end{equation}

The two quadratures affect LIGO in different ways:

\begin{enumerate}
    \item The phase quadrature \(\hat Y\) produces uncertainty in the
    optical phase measured at the output. This is called shot noise.

    \item The amplitude quadrature \(\hat X\) produces fluctuations in the
    number of photons incident on the mirrors. The corresponding fluctuating
    radiation pressure moves the mirrors. This is called quantum
    radiation-pressure noise.
\end{enumerate}

The approximate frequency dependence of the quadrature variance is
\begin{equation}
    S_h^{\text{shot}}(\Omega)\sim\Delta Y^2\propto\frac{1}{P},
    \qquad
    S_h^{\text{rad}}(\Omega)\sim\Delta X^2\propto\frac{P}{\Omega^4},
\end{equation}
where \(P\) is the circulating optical power and \(\Omega\) is the
gravitational-wave angular frequency. Shot noise is most important at high
frequencies, while radiation-pressure noise becomes important at low
frequencies.

Increasing the laser power reduces shot noise but increases
radiation-pressure noise. This tradeoff leads to the standard quantum limit.

\subsubsection{Injection of squeezed vacuum}

For ordinary vacuum,
\begin{equation}
    (\Delta X)^2=(\Delta Y)^2=\frac12.
\end{equation}
A squeezed vacuum state redistributes these uncertainties. For phase
squeezing,
\begin{equation}
    (\Delta X)^2=\frac12e^{2r},
    \qquad
    (\Delta Y)^2=\frac12e^{-2r},
\end{equation}
where \(r>0\) is the squeeze parameter. Therefore,
\begin{equation}
    \Delta X\,\Delta Y=\frac12.
\end{equation}

Injecting this squeezed vacuum into the antisymmetric port reduces the phase
uncertainty:
\begin{equation}
    \Delta Y_{\text{sq}}=e^{-r}\Delta Y_{\text{vac}}.
\end{equation}
The shot-noise contribution is consequently reduced by
\begin{equation}
    S_{h,\text{sq}}^{\text{shot}}
    =e^{-2r}S_{h,\text vac}^{\text{shot}}.
\end{equation}

The uncertainty principle is not violated. The reduction of phase noise is
accompanied by an increase in amplitude noise:
\begin{equation}
    \Delta X_{\text{sq}}=e^r\Delta X_{\text{vac}},~~~S_{h,\text{sq}}^{\text{rad}}=e^{2r}S_{h,\text{sq}}^{\text{rad}}.
\end{equation}
The increased amplitude noise can increase radiation-pressure noise at low
frequencies.

\subsubsection{Frequency-dependent squeezing}

A fixed squeezing angle cannot reduce both forms of quantum noise over the
entire detection band. LIGO therefore uses frequency-dependent squeezing.

Let the quadrature vector be
\begin{equation}
    \hat{\bm q}
    =
    \begin{pmatrix}
        \hat X\\
        \hat Y
    \end{pmatrix}.
\end{equation}
A rotation of the squeezing ellipse through an angle \(\theta\) is represented
by
\begin{equation}
    \hat{\bm q}_\theta
    =
    \begin{pmatrix}
        \cos\theta & -\sin\theta\\
        \sin\theta & \cos\theta
    \end{pmatrix}
    \hat{\bm q}.
\end{equation}

Before the squeezed vacuum enters the main interferometer, it is reflected
from an auxiliary optical filter cavity. The cavity produces a
frequency-dependent quadrature rotation,
\begin{equation}
    \theta=\theta(\Omega).
\end{equation}
The squeezing ellipse is therefore oriented differently for different
gravitational-wave frequencies.

At high frequencies, the squeezed quadrature is aligned to reduce phase
noise:
\begin{equation}
    (\Delta Y)^2_{\text{high\ frequency}}
    =\frac12e^{-2r}.
\end{equation}
This reduces shot noise.

At low frequencies, the filter cavity rotates the ellipse so that the
amplitude quadrature is squeezed:
\begin{equation}
    (\Delta X)^2_{\text{low\ frequency}}
    =\frac12e^{-2r}.
\end{equation}
This reduces radiation-pressure fluctuations acting on the mirrors.

Schematically,
\begin{equation}
    \boxed{
    \begin{aligned}
        \text{low frequency:}\quad
        &\text{amplitude squeezing}
        &&\Longrightarrow
        \text{less radiation-pressure noise},\\
        \text{high frequency:}\quad
        &\text{phase squeezing}
        &&\Longrightarrow
        \text{less shot noise}.
    \end{aligned}}
\end{equation}

\subsubsection{Optical loss}

Optical loss replaces part of the squeezed state with ordinary vacuum. If
\(\eta\) is the total optical efficiency, the observed squeezed variance is
approximately
\begin{equation}
    (\Delta Q)^2_{\text obs}
    =
    \eta\,\frac{e^{-2r}}{2}
    +(1-\eta)\frac12,
\end{equation}
where \(\hat Q\) denotes the intended squeezed quadrature.

Even if the generated squeezing is very strong, the measured noise reduction
is limited by optical loss and fluctuations of the squeezing angle. High
optical efficiency and accurate control of the squeezing phase are therefore
essential.

\subsubsection{Summary}

LIGO does not evade the uncertainty principle. Instead, squeezed light
places less quantum uncertainty in the quadrature that limits a particular
part of the measurement and more uncertainty in the conjugate quadrature.

Frequency-dependent squeezing extends this idea by rotating the squeezing
ellipse as a function of frequency. It reduces the relevant quantum noise at
both low and high frequencies, allowing weaker gravitational-wave signals to
be distinguished from the quantum fluctuations of the light.


This section summarizes LIGO’s use of vacuum fluctuations entering the unused interferometer port, squeezed states to reduce shot noise, and frequency-dependent squeezing to control the tradeoff between high-frequency phase noise and low-frequency radiation-pressure noise.

LIGO reports that frequency-dependent squeezing has been operating since the observing run began in May 2023 and improves sensitivity across its gravitational-wave frequency range.



